\subsection{位移}

我们再回忆，保持任意两点间距离不变的仿射变换 \(f\) （\(\mathbf{R}^n\) 到其自身的）{[}也就是说，使得 \(d(f(x), f(y)) = d(x, y)\) 对一切 \(x, y]\) 成立，称为欧氏位移（或简称位移）（*）；它们组成一个群，称为 \(\mathbf{R}^n\) 的位移群。这个群在 \(\mathbf{R}^n\) 上传递地作用；更一般地，若 \(V\) 与 \(V'\) 是任意两个同维的仿射线性簇（\(\mathbf{R}^n\) 中的），则存在把 \(V\) 变换为 \(V'\) 的一个位移。保持原点不动的位移称为正交变换，它们构成全体位移的群的一个子群。这个子群称为 \(n\) 个实变量的正交群；属于这个群的线性映射的特征是：它们保持范数 \(||x||\) ——每一点 \(x \in \mathbf{R}^n\) 的范数——不变，或者等价地，保持二次型 \(||x||^2 = \sum_{i=1}^{n} x_i^2\) 不变。两个向量 \(x = (x_i)\) 与 \(y = (y_i)\) （\(\mathbf{R}^n\) 的）的内积是值 \(\sum_{i=1}^{n} x_i y_i\) ，即与二次型 \(\frac{1}{2} \sum_{i=1}^{n} x_i^2\) 相伴的双线性形式所取的值；它记作 \((x|y)\) ，或在不会混淆时记作 \(xy\) 。每个正交变换保持任意两个向量的内积不变。两个向量 \(x, y\) 称为正交的，若 \((x|y) = 0\) ；两个向量子空间 \(V, V'\) （\(\mathbf{R}^n\) 的）称为正交的，若每个 \(x \in V\) 与每个 \(y \in V'\) 正交；两个仿射线性簇 \(P, P'\) 称为正交的，若分别与 \(P\) 和 \(P'\) 平行的向量子空间是正交的。

